\section{The broad pure-quartic reflection reduction on the same window}
\label{sec:quartic-abelian-reflection-actual-window}

\paragraph{Status and literal dependency.}
This conditional-interface proof applies the original-coefficient
strict theorem AC-G170, including its separate residue-zero argument.
It retains every omission of AC-G169-window and broadens the two
previous pure quartic rows to all four actual quartic labels.
The ten uncovered label--residue pairs are retained in the unknown sum.
One full local comparison cost covers the entire enlarged union.
The remaining all-prime estimate and mixed realization stay open.

Put $p=269$, $K=\mathbb Q(\omega)$ and $t=\sqrt[3]p>0$;
let $\chi$ describe $K(t)/K$.
Let $\phi$ be the primitive odd quartic character modulo sixteen
with $\phi(13)=i$ and kernel $\{1,7\}$, and set
\[
 \delta_d=\chi_d|_{G_K}\quad(d=1,-4,8,-8),\qquad
 \theta_{\sigma,a}=\phi^\sigma\chi_{-4}^{\,a}|_{G_K}
                    \quad(\sigma=\pm1,\ a=0,1).
 \tag{QRW.1}
\]
Here $\delta_1=1$.
For the retained datum, $0<\varepsilon<1/256$, page and
prescribed pair $q,d$, keep exactly
\[
 \begin{aligned}
 F_u&=\mathbb Q(\zeta_{2^{h(u)+2}},\zeta_7,
                                  \sqrt{-q},\sqrt{-d},\omega),\\
 H_u&=F_u(t),\quad G_u=\operatorname{Gal}(H_u/K),\quad D_u=|G_u|,\\
 u^{2+8\varepsilon}&\leq2^{h(u)}\leq2u^{2+8\varepsilon},
                    \qquad T=\lfloor u^{399}\rfloor .
 \end{aligned}
 \tag{QRW.2}
\]
The prescribed pair's $d$ is separate from the subscript of $\delta_d$.
All original $O_\ell,E_\ell,b_\ell$, full local conditions,
indices and unrestricted coefficient-prime range remain (FC.1)--(FC.2),
(EU.2) and (TC.1).

Define the already retained and the new exceptional residues by
\[
 \begin{aligned}
 \mathcal E&=\{89,92,99,168,175,178\},\\
 \mathcal F_{1,0}&=\{243,248\},&\mathcal F_{-1,0}&=\{19,24\},\\
 \mathcal F_{1,1}&=\{11,60,229\},&\mathcal F_{-1,1}&=\{38,207,256\}.
 \end{aligned}
 \tag{QRW.3}
\]
All congruences specifying these exceptions are modulo $268$.
For $u\geq2$ the complete retained part, with the old pure
quartic rows removed for replacement, is the following disjoint union:
\[
 \begin{aligned}
 Z_{16}^+&=\{(p,\chi^{-j}\phi^{-\sigma},e):j=1,2,\ \sigma=\pm1,\\
 &\hspace{2em}1\leq e<T,\ e\equiv-1\pmod{67}\},\\
 Z_{\rm quad}^+&=\{(p,\chi^{-j}\delta_d,e):j=1,2,\ d=-4,8,-8,\\
 &\hspace{2em}1\leq e<T,\ e\equiv-1\pmod{67}\},\\
 Z_{\rm cub}^{\ne0}&=\{(p,\chi^j,e):j=1,2,\ 1\leq e<T,\\
 &\hspace{2em}e\equiv0\pmod{268}\ \text{or}\ e\equiv-1\pmod{67}\},\\
 Z_1^{\rm all}&=\{(p,1,e):1\leq e<T\},\\
 Z_d^{\rm ref}&=\{(p,\delta_d,e):1\leq e<T,\ e\not\equiv0\pmod{268}\}
                                                \quad(d=-4,8),\\
 Z_{-8}^{\rm ref}&=\{(p,\delta_{-8},e):1\leq e<T,
                         \ e\bmod268\notin\{0\}\cup\mathcal E\},\\
 Z_3&=\{(3,1,e):1\leq e<T\},\\
 Z_{\rm retained}&=Z_{16}^+\sqcup Z_{\rm quad}^+
    \sqcup Z_{\rm cub}^{\ne0}\\
    &\hspace{2em}\sqcup Z_1^{\rm all}\sqcup Z_{-4}^{\rm ref}
       \sqcup Z_8^{\rm ref}\sqcup Z_{-8}^{\rm ref}\sqcup Z_3.
 \end{aligned}
 \tag{QRW.4}
\]
The replacement pure quartic set and full new union are
\[
 \begin{aligned}
 Z_{\rm quart}^{\rm ref}
   &=\{(p,\theta_{\sigma,a},e):\sigma=\pm1,\ a=0,1,\ 1\leq e<T,\\
   &\hspace{7em}e\bmod268\notin\mathcal F_{\sigma,a}\},\\
 Z^{\rm qref}&=Z_{\rm retained}\sqcup Z_{\rm quart}^{\rm ref}.
 \end{aligned}
 \tag{QRW.5}
\]
The predecessor $Z^{\rm ref}$ is (ARW.4) of AC-G169.
Its two pure odd quartic rows $Z_{16}^{0,q}$ have residues
$66,200$ only; all are in (QRW.5).
Hence $Z^{\rm ref}\subseteq Z^{\rm qref}$.
Let $I_{399}^{\rm estquartreflection}$ omit precisely $Z^{\rm qref}$
from (TC.1); its predecessor is $I_{399}^{\rm estreflection}$.

For positive integer $T$ put
\[
 \begin{aligned}
 N_T&=T-1,&B_T&=\lfloor(T-1)/268\rfloor,\\
 A_T&=\lfloor T/67\rfloor,&Q_T&=\lfloor(T+67)/134\rfloor,\\
 R_T&=\sum_{r\in\mathcal E}\left\lfloor\frac{T+267-r}{268}\right\rfloor,\\
 S_T&=\sum_{\sigma=\pm1}\sum_{a=0}^1
           \sum_{r\in\mathcal F_{\sigma,a}}
                    \left\lfloor\frac{T+267-r}{268}\right\rfloor.
 \end{aligned}
 \tag{QRW.6}
\]
Here $S_T$ counts exceptional label--weight pairs, not merely
unlabelled residue classes.

\begin{theorem}[AC-G171-window: the broad pure-quartic reflection reduction]
\label{AC-G171-window}
Every triple in (QRW.4)--(QRW.5) is an actual host character,
coefficient prime and weight with zero original strict group.
Every AC-G169 omitted term is retained.
The exact new pure quartic, new relative and cumulative counts are
\[
 \begin{aligned}
 |Z_{\rm quart}^{\rm ref}|&=4N_T-S_T,\\
 |Z^{\rm qref}\setminus Z^{\rm ref}|&=4N_T-S_T-2Q_T,\\
 |Z^{\rm qref}_{269}|&=8N_T-B_T+12A_T-R_T-S_T,\\
 |Z^{\rm qref}_3|&=N_T.
 \end{aligned}
 \tag{QRW.7}
\]
The inclusion is proper for $T\geq2$.
No zero or nonzero assertion is made at the ten uncovered
pure quartic pairs or any other remaining term.

Let $C_u$ be the actual normalized positive-index cost over
the entire union $Z^{\rm qref}$.  In nonnegative extended sums,
\[
 \begin{aligned}
 I_{399}&=I_{399}^{\rm estquartreflection}+C_u,\\
 I_{399}^{\rm estquartreflection}&\leq I_{399}^{\rm estreflection}\leq I_{399},\\
 0\leq C_u&\leq8D_u^2T^2\log(8TD_u)
                   +\tfrac12D_u\mathcal R_uT^2+2D_ur_uc_0T\\
 &=O_\varepsilon\bigl(u^{802+16\varepsilon}(\log u)^2\bigr).
 \end{aligned}
 \tag{QRW.8}
\]
This is one full (FL.4) error for the entire old and new union;
no individual positive index is asserted zero.
The new positive-relative-slack premise implies the full original
common-window premise on the same $W,J$, datum, page and pair,
after increasing only the initial threshold and halving slack.
Every older premise implies the new one by nonnegativity.
All simultaneous field, original coefficient, finite-character and
weight quantifiers, the actual $R.t$ and the retained mixed input
$\mathcal R_{\rm cub}$ remain unchanged.
\end{theorem}

\begin{proof}
\proofstep{1. Actual labels and literal original-coefficient strict zero.}
For $u\geq2$, (QRW.2) gives $h(u)\geq2$, so the actual host
contains $\mathbb Q(\zeta_{16})$.
Its four primitive quartic characters are exactly
$\phi^\sigma\chi_{-4}^{\,a}$ with $\sigma=\pm1$ and $a=0,1$.
They remain distinct over $K$, since
$K\cap\mathbb Q(\zeta_{16})=\mathbb Q$.
They differ from every quadratic or trivial label.
Also $F_u\cap K(t)=K$: the rational abelian field $F_u$
cannot contain the nonabelian splitting field $K(t)/\mathbb Q$.
Thus they differ from every retained nontrivial-cubic label.
The coefficient-three row is separate.
These facts prove the disjointness in (QRW.4)--(QRW.5).

For a pure quartic label the actual coefficient is
\[
 \begin{aligned}
 A_{\sigma,a,e}
   &=(O_p\otimes_{\mathbb Z_p}\mathbb Q_p/\mathbb Z_p)
                                           (\theta_{\sigma,a})(-e)\\
   &=(E_p/O_p)(\theta_{\sigma,a}c_p^{-e}).
 \end{aligned}
 \tag{QRW.9}
\]
Apply (G170.4) with precisely $(\sigma,a)$, $m=0$ and $\kappa=1$.
It proves the strict zero of (QRW.9) for every weight in
$Z_{\rm quart}^{\rm ref}$, on this same original $O_p$ with
arbitrary absolute ramification.
At nonzero residue its proof uses residual cyclotomic inertia;
at residue zero it separately uses
$\theta_{\sigma,a}(\operatorname{Frob}_{K_p})=-1$.
Thus the residue-zero assertion is original-coefficient and is
not inferred merely from a residual class calculation.
The auxiliary field $\mathbb Q(\zeta_{48p})$ is not substituted
for the actual host, and no new conductor-$p$ finite label is added.
Every strict zero in $Z_{\rm retained}$ is kept from AC-G169,
including its separate old trivial-label residue-zero input.

The source imposes zero in full $H^1(G_{K_v},A)$ at $v\mid p$
and zero in $H^1(I_v,A)$ at every other finite place.
These are the actual (FC.2) conditions.
For added away-$p$ ramification primes the coefficient is
unramified; a cohomologically zero inertia restriction is then
a zero cocycle.  The cocycle identity gives zero on conjugate
inertia groups and their closed normal closure, so the class
descends to the smaller ramification set.
Conversely inflation is injective and preserves full local zero.
This proves the required equality of strict kernels on the
original coefficient, including the growing host's enlarged set.

\proofstep{2. Exact replacement and counts.}
All ten exceptional residues are nonzero.
For each original quartic label the count of retained weights is
$N_T-\sum_{r\in\mathcal F_{\sigma,a}}
\lfloor(T+267-r)/268\rfloor$, by counting positive integers
$e<T$ in each missing progression.  Summing gives $4N_T-S_T$.
For $a=0$ none of the four exceptions $19,24,243,248$
equals $66$ or $200$.
Consequently all old pure quartic rows lie in the replacement,
with old count $2Q_T$.
Subtracting these finite set cardinalities gives the second line
of (QRW.7), with no double counting.
The old coefficient-$269$ count from (ARW.7) was
$4N_T-B_T+12A_T+2Q_T-R_T$.
Replacing its $2Q_T$ contribution by $4N_T-S_T$ gives the
third line of (QRW.7); coefficient three is unchanged.
At $e=1$ all four pure quartic labels are new, proving proper
inclusion for $T\geq2$.  This need not strictly decrease a
nonnegative numerical sum if its added terms already vanish.

\proofstep{3. One full-union comparison cost.}
Retain the codifferent identification and original normalizing
rank $b_\ell$ in the dual complex (FL.2).
For each omitted triple, its literal strict zero and (EU.2) imply
\[
 \begin{aligned}
 \frac{\log\#k_{O_\ell}}{b_\ell}\max(0,i_{\ell,\eta,e})
 &\leq\frac{\log\#H^2(C_{\ell,\eta,e})}{b_\ell}\\
 &=\frac{\left|\log\#H^2(C_{\ell,\eta,e})
                  -\log\#\mathcal S_{\ell,\eta}(e)\right|}{b_\ell}.
 \end{aligned}
 \tag{QRW.10}
\]
At each weight the entire union is one subsum of the nonnegative
all-prime (FL.4) error.  Its bound is
$16eD_u^2\log(8eD_u)+D_ue\mathcal R_u+2D_ur_uc_0$.
Summing once over $1\leq e<T$ gives (QRW.8), using
$\sum e\leq T^2/2$.
No new full-error allowance is added for the quartic replacement.
The extended-sum identity and inequalities follow by nonnegative
additivity, without subtracting infinities.
The unchanged $D_u\leq288u^{2+8\varepsilon}$, $r_u\leq12$,
$\mathcal R_u=O_\varepsilon(u^{\varepsilon/5}(\log u)^2)$ and
$T\leq u^{399}$ give the displayed exponent $802+16\varepsilon$.
This cost is $o(\exp(u^\varepsilon/2))$ for each fixed positive
$\varepsilon$, uniformly in the prescribed pair.

\proofstep{4. The unchanged common window.}
Fix the retained datum, page and $0<\varepsilon<1/256$.
Suppose the new premise supplies $J\geq1$, $X_0$ and
$s\in(0,1)$, independent of $X,q,d,u,\ell,\eta,e$, such that
for every $X\geq X_0$ and prescribed
$\operatorname{WinPair}(\varepsilon,X^{1/J},q,d)$ there is
one $W$ with $X^{1/J}\leq W$, $W^2\leq X$ and
\[
 I_{399}^{\rm estquartreflection}(H_u,u)
       \leq(1-s)\exp(u^\varepsilon/2)\qquad(W\leq u\leq W^2).
 \tag{QRW.11}
\]
Choose $U_0\geq2$ so that (QRW.8)'s cost is at most
$(s/2)\exp(u^\varepsilon/2)$ for every $u\geq U_0$.
Replacing $X_0$ by $\max(X_0,U_0^J)$ leaves the same $W$
and proves the original full bound with slack $s/2$ there.
For a fixed finite simultaneous family, take the maximum of
its uniform error constants on that same window.
No intersection of independently selected windows is used.
Conversely either older premise implies the new one by
nonnegativity, with unchanged parameters.
The same reasoning covers uniform little-$o$ premises.
All actual fields, original coefficients, source $R.t$ and
mixed-realization clauses remain exactly as before.
\end{proof}

\paragraph{Remaining scope and geometric dictionary.}
The ten unsupported pure quartic label--residue pairs remain,
as do every unlisted coefficient prime, character and weight,
the three nontrivial pure quadratic residue-zero families and
the six unsupported $\delta_{-8}$ residues.
This does not assert nonvanishing or a globally maximal classification.
The remaining numerical upper bound, mixed theta/Selmer identification,
analytic realization and H04.4 remain open.
The previously audited SGA~1, Expos\'e~V, \S\S4 and 7 finite-cover
and original rank-one Artin--Tate dictionary is retained.
No new motive, cohomology comparison or Lean execution is asserted.
